\documentclass[hidelinks]{article}
\usepackage[a4paper, total={7in, 10in}]{geometry}
\usepackage[dvipsnames]{xcolor}
\usepackage{amsmath, amssymb, amsthm}
\usepackage{tikz}
\usepackage{tkz-euclide}
\usepackage[unicode]{hyperref}
\usepackage[all]{hypcap}
\usepackage{fancyhdr}
\usepackage{listings}
\usepackage{subcaption}
\usepackage{enumitem}
\usepackage{graphicx}

\lstset{basicstyle=\ttfamily\small,columns=fullflexible,frame=single,breaklines=true,showstringspaces=false}

\title{\textbf{Electromagnetic Radiation Test}}
\author{Diego Juarez}
\date{March 25th, 2026}

\pagestyle{fancy}
\fancyhf{}
\fancyhead[L]{Electromagnetic Radiation Test}
\fancyhead[R]{Jackson-style Practice}
\fancyfoot[C]{\thepage}

\begin{document}
\hypersetup{bookmarksnumbered=true}
\pagecolor{black}
\color{white}
\maketitle

\begin{Large}
\tableofcontents
\end{Large}
\pagebreak

\section{Source Notebook}
This problem set follows the notebook
\url{../notebooks/Electromagnetic Radiation.ipynb}
and focuses on retarded potentials, radiation-zone fields, dipole radiation, and power flow.

\section{Foundations}
\subsection{Problem 1: Wave Equations from Maxwell's Equations}
Starting from Maxwell's equations in vacuum, derive the wave equations for $\mathbf{E}$ and $\mathbf{B}$ in source-free space.

\subsection{Problem 2: Retarded Potentials}
Write the scalar and vector retarded potentials in the Lorenz gauge for general charge and current densities. Explain the meaning of retarded time.

\subsection{Problem 3: Causality}
Why are retarded rather than advanced potentials used in classical electrodynamics? Give a physical explanation rather than a formal one only.

\section{Radiation Zone}
\subsection{Problem 4: Near, Induction, and Radiation Terms}
For a localized oscillating source, explain the qualitative difference between terms that scale as
\[
\frac{1}{r^3},\qquad \frac{1}{r^2},\qquad \frac{1}{r}.
\]
Which term carries energy to infinity?

\subsection{Problem 5: Far-Field Approximation}
State the approximations that define the radiation zone. Explain why in that region the electric and magnetic fields are transverse.

\subsection{Problem 6: Poynting Vector}
Starting from the far-zone relation between $\mathbf{E}$ and $\mathbf{B}$, show that the Poynting vector points radially outward on average for outgoing radiation.

\section{Dipole Radiation}
\subsection{Problem 7: Oscillating Electric Dipole}
An electric dipole moment oscillates as
\[
\mathbf{p}(t)=p_0\cos(\omega t)\,\hat{\mathbf{z}}.
\]
\begin{enumerate}[label=\alph*)]
\item Write the far-zone angular dependence of the radiated electric field.
\item Write the corresponding magnetic field.
\item State the dependence on $r$, $\omega$, and $\sin\theta$.
\end{enumerate}

\subsection{Problem 8: Radiation Pattern}
For the oscillating dipole oriented along the $z$-axis:
\begin{enumerate}[label=\alph*)]
\item At which angles is the radiation maximal?
\item At which angles does it vanish?
\item Sketch the qualitative angular pattern.
\end{enumerate}

\subsection{Problem 9: Total Radiated Power}
Use the angular distribution of dipole radiation to derive the total time-averaged radiated power of a nonrelativistic oscillating dipole.

\subsection{Problem 10: Hertzian Dipole}
Explain what is meant by a Hertzian dipole. Why is it a useful elementary source in radiation theory?

\section{Physical Interpretation}
\subsection{Problem 11: Accelerating Charge}
Explain why a uniformly moving charge does not radiate, while an accelerating charge can. Your answer should distinguish clearly between static, induction, and radiation behavior.

\subsection{Problem 12: Multipole Viewpoint}
Why is electric dipole radiation usually the leading radiation mechanism for a localized oscillating source? Under what symmetry conditions can dipole radiation vanish?

\subsection{Problem 13: Radiation Reaction Discussion}
What is meant by radiation reaction? Explain why it is conceptually subtle even before writing down a detailed formula.

\section{Computation-Oriented Problems}
\subsection{Problem 14: Frequency Dependence}
How does the total radiated power of an oscillating dipole scale with angular frequency $\omega$ and dipole amplitude $p_0$? What does this imply about high-frequency sources?

\subsection{Problem 15: Near-Field Versus Far-Field}
Suppose measurements are made at several distances from a source. Explain how one could distinguish experimentally between near-field and far-field behavior using radial scaling.

\section{Challenge Problems}
\subsection{Problem 16: Symmetry and Forbidden Dipole Radiation}
Give one example of a source configuration for which the electric dipole moment vanishes by symmetry. What multipole order would then typically dominate the radiation?

\subsection{Problem 17: Dimensional Analysis of Larmor-Type Power}
Use dimensional analysis and basic physical reasoning to motivate the dependence of radiated power on charge, acceleration, vacuum permittivity, and the speed of light for a nonrelativistic point charge.

\end{document}
